\chapter{Geometric Group Theory}\label{geometric-group-theory}
\chapterstatus{planned}

Finitely generated groups studied as geometric objects: through their Cayley graphs, their actions on spaces, and invariants of their large-scale geometry, such as growth, hyperbolicity, ends and
amenability.

\noindent\textbf{Prerequisites.} Chapter~\ref{groups} (Groups), Chapter~\ref{homotopy} (Homotopy and Covering Spaces), Chapter~\ref{classical-geometry} (Classical Geometry) and
Chapter~\ref{low-dimensional-topology} (Low-Dimensional Topology).

\noindent\textbf{Planned contents.}
\begin{enumerate}
\item Cayley graphs and word metrics; quasi-isometries; the Švarc--Milnor lemma; quasi-isometry invariants.
\item Groups acting on trees: Bass--Serre theory, amalgamated products and HNN extensions, $\mathrm{SL}_2(\ZZ) \cong \ZZ/4 *_{\ZZ/2} \ZZ/6$, free groups and Nielsen--Schreier revisited.
\item Growth: polynomial and exponential growth, Milnor--Wolf, Gromov's polynomial growth theorem (statement), Grigorchuk's group of intermediate growth.
\item Hyperbolic groups: Gromov hyperbolic spaces, the Morse lemma, the word and conjugacy problems, Dehn's algorithm, small cancellation theory, the boundary.
\item Ends of groups and Stallings' theorem (statement).
\item Amenability: Følner sets, the von Neumann--Day problem, paradoxical decompositions and the Banach--Tarski theorem revisited, Tits alternative (statement), Kazhdan's property (T) (survey).
\item Exercises with solutions.
\end{enumerate}
